\documentclass[10pt,a4paper]{article}
\usepackage{nexusS5}
\setnexusheader{Entrée en Première générale — Spécialité mathématiques — Stage de pré-rentrée}{Évaluation finale — Malek KHADHRANI}
\begin{document}
\nexuscover{ÉVALUATION FINALE --- DIAGNOSTIC DE SORTIE}{Durée : 45 minutes}{Malek KHADHRANI}{Entrée en Première générale — Spécialité mathématiques}{Mathématiques --- à traiter individuellement}
\par\noindent\begin{tabularx}{\linewidth}{>{\bfseries}p{22mm} X >{\bfseries}p{16mm} X}
\toprule Nom et prénom & \rule{62mm}{0.32pt} & Date & \rule{34mm}{0.32pt} \\ \bottomrule
\end{tabularx}

\begin{goalbox}
\textbf{Ce document est un diagnostic de sortie, pas un classement.} Il mesure ce qui est
désormais installé, ce qui reste fragile, et il sert à écrire le plan des quatre premières
semaines de la rentrée.
\end{goalbox}

\begin{methodbox}
\textbf{Consignes}
\begin{itemize}
\item Durée : \textbf{45 minutes}. Partie A : environ 10 min. Partie B : environ 20 min. Partie C : environ 15 min.
\item Travailler seul, sans document et sans les cartes de rappel de la séance.
\item Écrire la relation ou la propriété utilisée avant le calcul.
\item Une question peut être laissée de côté puis reprise ; il vaut mieux la laisser vide que d'écrire au hasard.
\item Trois lignes « Certitude » seulement : \conf\ , sur la même échelle qu'en début de stage.
\end{itemize}
\end{methodbox}

\newpage
\phasehead{Partie A --- Automatismes et connaissances essentielles}{environ 10 min}{Questions courtes. Écrire le résultat, sans rédaction longue.}
\begin{enumerate}
\needspace{22mm}
\item[\textbf{A1.}] Factoriser $16x^2 - 49$.
\worklines{2}
\par\vspace{1.2mm}
\needspace{22mm}
\item[\textbf{A2.}] Résoudre $-3x + 12 > 0$ et donner l'ensemble des solutions sous forme d'intervalle.
\worklines{2}
\par\vspace{1.2mm}
\needspace{22mm}
\item[\textbf{A3.}] Soit $g(x) = 2x^2 + x$. Calculer $g(-3)$ en écrivant la substitution complète.
\worklines{2}
\par\vspace{1.2mm}
\needspace{22mm}
\item[\textbf{A4.}] On donne $A(4\,;-1)$ et $B(-2\,;5)$. Déterminer les coordonnées de $\vec{AB}$.
\worklines{2}
\par\vspace{1.2mm}
\needspace{22mm}
\item[\textbf{A5.}] Soit $f(x) = x^2 - 6x$. Calculer $f(-4)$ en écrivant la substitution complète.
\worklines{2}
\par\vspace{1.2mm}
\needspace{22mm}
\item[\textbf{A6.}] Les vecteurs $\vec{u}(6\,;-4)$ et $\vec{v}(-9\,;6)$ sont-ils colinéaires ? Justifier par le calcul du déterminant.
\worklines{2}
\par\vspace{1.2mm}
\end{enumerate}
\certbox{Certitude sur l'ensemble de la partie A :}
\needspace{68mm}\par\vspace{3mm}
\phasehead{Partie B --- Application et raisonnement}{environ 20 min}{Écrire la relation, la propriété ou la règle avant le calcul, puis contrôler.}
\begin{enumerate}
\needspace{22mm}
\item[\textbf{B1.}] On considère le produit $P(x) = (3x + 6)(2 - x)$.
\begin{enumerate}
\item Déterminer les valeurs de $x$ qui annulent chaque facteur.
\item Construire un tableau de signes, puis donner l'ensemble des $x$ pour lesquels $P(x) > 0$.
\item En déduire, à partir du même tableau, l'ensemble des $x$ pour lesquels $P(x) \leqslant 0$.
\end{enumerate}
\worklines{6}
\par\vspace{1.2mm}
\needspace{22mm}
\item[\textbf{B2.}] \begin{enumerate}
\item La courbe représentative d'une fonction $f$ passe par le point de coordonnées $(-2\,;5)$.
Traduire cette information en employant les mots « image » et « antécédent ».
\item Soit $x$ un réel tel que $0 < x < 1$. Comparer $x^2$ et $x$, puis justifier la réponse
par un exemple numérique et par un argument valable pour tout $x$ de cet intervalle.
\end{enumerate}
\worklines{6}
\par\vspace{1.2mm}
\needspace{22mm}
\item[\textbf{B3.}] Un abonnement coûte $480$ TND. Son prix baisse de $25\,\%$, puis le nouveau prix
augmente de $20\,\%$.
\begin{enumerate}
\item Calculer le prix après chaque étape.
\item Donner le coefficient multiplicateur global, puis le taux d'évolution global en pourcentage.
\end{enumerate}
\worklines{5}
\par\vspace{1.2mm}
\needspace{22mm}
\item[\textbf{B4.}] On donne $A(3\,;-2)$, $B(7\,;6)$ et $C(1\,;-6)$.
\begin{enumerate}
\item Calculer les coordonnées de $\vec{AB}$ et de $\vec{AC}$.
\item Les points $A$, $B$ et $C$ sont-ils alignés ? Justifier.
\item Donner le coefficient directeur de la droite $(AB)$.
\end{enumerate}
\worklines{7}
\par\vspace{1.2mm}
\end{enumerate}
\certbox{Certitude sur la question B2 :}
\needspace{68mm}\par\vspace{3mm}
\phasehead{Partie C --- Transfert et synthèse}{environ 15 min}{Reconnaître la notion utile, mobiliser plusieurs acquis, conclure par une phrase.}
\begin{enumerate}
\needspace{22mm}
\item[\textbf{C1.}] \textbf{Le club sportif.} Dans un repère du plan, on donne $A(1\,;2)$ et $B(5\,;10)$.
\begin{enumerate}
\item Déterminer le coefficient directeur de la droite $(AB)$, puis une équation réduite de $(AB)$.
\item Le point $C(-3\,;-6)$ appartient-il à $(AB)$ ? Justifier par un calcul, en indiquant la méthode choisie.
\item Le club compte $60$ membres : $24$ pratiquent la natation, $30$ la course, et $12$ pratiquent les deux.
On choisit un membre au hasard. Calculer la probabilité qu'il pratique la natation ou la course,
puis la probabilité qu'il ne pratique pas la natation.
\item La cotisation, de $180$ TND, augmente de $4\,\%$ par an. On note $u_0 = 180$. Écrire la relation liant
$u_{n+1}$ et $u_n$, puis calculer $u_2$ arrondi au dinar.
\end{enumerate}
\worklines{12}
\par\vspace{1.2mm}
\needspace{22mm}
\item[\textbf{C2.}] \begin{enumerate}
\item Simplifier $\dfrac{3^{5}\times 3^{-2}}{3^{4}}$ et donner le résultat sous forme d'une puissance de $3$.
\item L'affirmation « la fonction inverse $f(x) = \dfrac{1}{x}$ est croissante sur $\left]0\,;+\infty\right[$ »
est-elle exacte ? Répondre par oui ou non, en justifiant à l'aide de deux valeurs bien choisies.
\end{enumerate}
\worklines{5}
\par\vspace{1.2mm}
\end{enumerate}
\certbox{Certitude sur la question C1 :}
\brouillon{\dimexpr\textheight/3\relax}
\newpage
\sectiontitle{Après l'évaluation --- à remplir en deux minutes}
\begin{methodbox}Ces trois lignes ne sont pas notées. Elles servent à construire le plan des quatre premières semaines de la rentrée.\end{methodbox}
\par\noindent\begin{tabularx}{\linewidth}{>{\bfseries}p{62mm} X}
\toprule
La question que j'ai trouvée la plus sûre & \rule{85mm}{0.32pt} \\
La question sur laquelle j'ai le plus hésité & \rule{85mm}{0.32pt} \\
Une notion que je veux revoir dès la première semaine & \rule{85mm}{0.32pt} \\
\bottomrule\end{tabularx}
\vfill
\begin{graybox}\small Ce document est un diagnostic de sortie. Il sera analysé compétence par compétence, puis comparé au positionnement passé en début de stage lorsque cette comparaison est légitime.\end{graybox}
\end{document}
